\section{Conclusion} \label{sec:conclusion}
This paper studies a model \`a la \citet{kyle1985continuous} where traders are privately informed about their skill, i.e., the trader-specific ability to identify potentially mispriced assets. The traders care not only about trading profits but also the recruiter's perception of their skill. The model has two stable self-fulfilling equilibria, which we interpret as fragility in financial markets. One of the equilibria is similar to that of \citet{kyle1985continuous} where reputation is irrelevant. In the other, ``reputation-driven'' equilibrium, the traders trade more aggressively for fear of being seen as unskilled by the recruiter, causing higher price volatility, higher market liquidity, and lower trading profits. Our model predicts that fragility may arise when poaching is (moderately) active in the finance industry and when there are high demands for financial skill. The model yields implications connecting traders’ skills, their
trading styles, and financial fragility.   Moreover, consistent with empirical observations, our model predicts large pay inequality in the finance industry and reconciles skill in active management and the lack of performance persistence.

As a direction for future research, it would be interesting to analyze various types of wage structures in financial markets, ranging from back-end to front-end compensations.  The implications of the forward-looking trading motives introduced in our model can be applied to investigate which types of wage structures and employment systems can mitigate financial market fragility while also improving the welfare of market participants.
